% Exact conjugacy-orbit resolution of the local C3 Clifford obstruction.
\subsection{The $C_9$ obstruction is one affine Clifford conjugacy orbit}

The local lift-charge theorem leaves a $16+24$ split among the $40$ cyclic order-three subgroups of
\[
Q\cong\operatorname{ASL}(2,3)=\mathbb F_3^2\rtimes\operatorname{SL}(2,3).
\]
Constructing this affine group directly resolves those $40$ subgroups into three conjugacy orbits,
\[
\boxed{40=4+12+24}.
\]
The orbit of size $4$ consists of the pure translations; their nonidentity elements have no affine fixed points.  The orbit of size $12$ consists of nontranslation unipotents with an affine fixed line, hence three fixed points.  The orbit of size $24$ consists of nontranslation unipotents with no affine fixed point.  The explicit quotient model has element-order distribution
\[
1^1\,2^9\,3^{80}\,4^{54}\,6^{72},
\]
matching the certified $216$-element Clifford quotient census.

Splitting of a central extension over a cyclic subgroup is invariant under quotient conjugacy.  Since the exact lift audit contains precisely $24$ nonsplit cyclic $C_3$ restrictions, the unique $24$-subgroup orbit must be the entire nonsplit locus.  Therefore
\[
\boxed{4_{\rm translation}+12_{\rm fixed\ line}\ \longmapsto\ C_3\times C_3}
\]
whereas
\[
\boxed{24_{\rm fixed\text{-}point\text{-}free\ nontranslation}\ \longmapsto\ C_9}.
\]
Thus the $3/5$ obstruction is geometrically homogeneous: the central deck charge is nontrivial exactly on the fixed-point-free order-three affine directions whose linear part is itself nontrivial unipotent.  Pure translations are also fixed-point-free, but they split, so fixed-point-freeness alone is not the discriminant.  This classification is exact in the canonical $\operatorname{ASL}(2,3)$ quotient model; an optical or hardware phase interpretation still requires a separately proved physical coordinate intertwiner.
